\documentclass[12pt]{article}

\usepackage[margin=1in]{geometry}
\usepackage{enumitem}
\usepackage{amsmath}
\usepackage{amssymb}
\usepackage{array}
\usepackage{titlesec}

\title{Lecture 11 Notes:\\ Aristotle's \textit{Topics}, Book IV}
\author{}
\date{}

\begin{document}

\maketitle

\section*{Central Theme}

Aristotle's \textit{Topics}, Book IV gives commonplace rules for testing
proposed genera.

\[
\boxed{
\textit{Topics}, \text{ Book IV}
=
\text{commonplace rules for testing proposed genera}
}
\]

Book IV teaches how to ask whether someone has placed a thing in the right
kind.

A genus must be wider than the species, predicated in essence, in the same
category, literally and synonymously said, accompanied by its higher genera, and
not confused with accident, differentia, capacity, activity, metaphor, part, or
mere accompaniment.

\section*{1. From Comparison to Genus}

Book III taught rules for comparing goods.

\[
\text{Is } A \text{ better or more desirable than } B?
\]

Book IV returns to the predicables and asks about genus.

\[
\text{Is } G \text{ really the genus of } S?
\]

\subsection*{Whiteboard}

\[
\text{Book III: comparison of goods}
\]

\[
\text{Book IV: testing proposed genera}
\]

\subsection*{Teaching Point}

To test a genus is to test whether a thing has been placed in the right kind.

\[
\boxed{
\text{A false genus produces false classification.}
}
\]

\section*{2. Why Genus Matters}

A genus tells us what something is broadly.

\[
\text{species} \subset \text{genus}
\]

\[
\text{man} \subset \text{animal}
\]

The genus does not give the complete definition, but it gives the higher kind
under which the species falls.

\subsection*{Teaching Point}

A genus is predicated in the category of essence.

\[
\boxed{
\text{Genus says what a thing is broadly.}
}
\]

\section*{3. The Basic Genus Question}

The basic question of Book IV is:

\[
\boxed{
\text{Is this proposed genus truly predicated of the species in essence?}
}
\]

That question can fail in many ways.

A proposed genus may really be:

\begin{itemize}[itemsep=0.25em]
    \item an accident,
    \item a quality,
    \item a relation,
    \item a metaphor,
    \item a differentia,
    \item a part,
    \item a capacity,
    \item an activity,
    \item a mere accompaniment.
\end{itemize}

\subsection*{Teaching Point}

Book IV trains us to detect false classification.

\[
\boxed{
\text{Do not confuse what a thing is with what happens to belong to it.}
}
\]

\section*{4. First Test: Is the Genus Predicated of Every Member?}

If a proposed genus does not apply to every member of the species, it is not the
genus.

\[
G \text{ is genus of } S
\Rightarrow
G \text{ belongs to every } S
\]

\[
\text{some } S \text{ is not } G
\Rightarrow
G \text{ is not the genus}
\]

\subsection*{Example}

Suppose someone says:

\[
\text{good is the genus of pleasure.}
\]

Then ask:

\[
\text{Is some pleasure not good?}
\]

If some pleasure is not good, then good is not the genus of pleasure.

\subsection*{Teaching Point}

One species-member outside the proposed genus refutes the genus.

\[
\boxed{
\text{A true genus is predicated of all members of the species.}
}
\]

\section*{5. Second Test: Essence or Accident?}

A proposed genus must be predicated in the category of essence, not merely as an
accident.

\subsection*{Example: Snow and White}

\[
\text{snow is white}
\]

But:

\[
\text{snow is not a kind of white}
\]

White is a quality of snow, not the genus of snow.

\subsection*{Example: Soul and Self-Moved}

\[
\text{soul may be self-moved}
\]

But:

\[
\text{soul is not a kind of moving object}
\]

Motion is not the essence of soul.

\subsection*{Teaching Point}

A genus says what a thing is, not merely what happens to belong to it.

\[
\boxed{
\text{Accidental predication is not generic predication.}
}
\]

\section*{6. Use the Definition of Accident}

An accident is what can belong or not belong to the same thing.

\[
\text{Accident} = \text{can belong or not belong}
\]

A genus does not come and go while the species remains.

\[
\text{Genus} = \text{belongs as long as the species exists}
\]

\subsection*{Examples}

A thing may be white or not white.

A thing may be self-moved or not self-moved.

Therefore, whiteness and self-motion are accidents, not genera, in these cases.

\subsection*{Teaching Point}

What can come and go cannot be the genus.

\[
\boxed{
\text{A genus belongs to the species essentially, not intermittently.}
}
\]

\section*{7. Genus and Species Must Be in the Same Category}

The genus and species must fall under the same categorical division.

If the species is a substance, the genus must be a substance.

If the species is a quality, the genus must be a quality.

If the species is a relative, the genus should be a relative.

\subsection*{Examples}

Snow and swan are substances.

White is a quality.

Therefore:

\[
\text{white is not the genus of snow or swan.}
\]

Knowledge is a relative.

Good and noble are qualities.

Therefore:

\[
\text{good and noble are not the genus of knowledge.}
\]

\subsection*{Teaching Point}

A genus must live in the same categorical division as its species.

\[
\boxed{
\text{Category mismatch refutes a proposed genus.}
}
\]

\section*{8. Species Partake of Genera}

To partake of something is to admit its definition.

A species partakes of its genus because the definition of the genus applies to
the species.

\[
\text{species partakes of genus}
\]

But the genus does not partake of the species.

\[
\text{genus does not partake of species}
\]

\subsection*{Example}

Man admits the definition of animal.

\[
\text{man} \rightarrow \text{animal}
\]

But animal does not admit the definition of man.

\[
\text{animal} \not\rightarrow \text{man}
\]

\subsection*{Teaching Point}

The genus is wider than the species.

\[
\boxed{
\text{Partaking runs upward from species to genus, not downward from genus to species.}
}
\]

\section*{9. The Genus Must Be Wider Than the Species}

The genus has wider denotation than the species and the differentia.

\[
\text{genus} > \text{species}
\]

If the proposed genus has equal or narrower extension, it is not the genus.

\[
\text{equal extension} \Rightarrow \text{not genus/species}
\]

\[
\text{narrower extension} \Rightarrow \text{not genus}
\]

\subsection*{Example}

Being and Unity are predicated of everything.

So neither can be the genus of the other in the ordinary way.

\[
\text{Being} = \text{Unity in extension}
\]

\[
\therefore \text{neither is the genus of the other}
\]

\subsection*{Teaching Point}

If the proposed genus is not wider than the species, it is not a genus.

\[
\boxed{
\text{Wider extension is a mark of genus.}
}
\]

\section*{10. Species Wider Than Genus Reverses the Order}

If the species has wider denotation than the proposed genus, the classification
has been reversed.

\subsection*{Example}

Object of opinion is wider than being, because both what exists and what does
not exist may be objects of opinion.

\[
\text{object of opinion} > \text{being}
\]

Therefore, being cannot be the genus of object of opinion.

\subsection*{Teaching Point}

The genus must include the species, not be included by it.

\[
\boxed{
\text{A narrower term cannot be the genus of a wider one.}
}
\]

\section*{11. If the Genus Applies, Some Species of the Genus Must Apply}

If a thing is placed under a genus, it must fall under one of the species
generated by the first division of that genus.

\[
S \subset G
\Rightarrow
S \subset \text{some species of } G
\]

\subsection*{Example}

Suppose someone says:

\[
\text{motion is the genus of pleasure.}
\]

Then pleasure must be some species of motion.

\[
\text{pleasure} =
\text{locomotion, alteration, growth, diminution, or another species of motion}
\]

If pleasure is none of these, then motion is not its genus.

\subsection*{Teaching Point}

A generic claim must be cashed out in the genus's proper divisions.

\[
\boxed{
\text{If no species of the genus applies, the genus does not apply.}
}
\]

\section*{12. Things Not Specifically Different Should Have the Same Genus}

All things that are not specifically different have the same genus.

\[
\text{not specifically different} \Rightarrow \text{same genus}
\]

If a proposed genus applies to one but not another thing not specifically
different, the genus claim fails.

\subsection*{Teaching Point}

A genus cannot arbitrarily split things that are not specifically different.

\[
\boxed{
\text{Sameness of species requires sameness of genus.}
}
\]

\section*{13. Multiple Genera Must Be Ordered}

If one species falls under two genera, then the genera should be ordered.

Either:

\[
G_1 \subset G_2
\]

or:

\[
G_2 \subset G_1
\]

or both fall under a higher genus:

\[
G_1, G_2 \subset H
\]

\subsection*{Example}

If justice is said to be both knowledge and virtue, then we must ask how
knowledge and virtue are related.

Are they subordinate to each other?

Do they both fall under some higher genus, such as state or disposition?

\subsection*{Teaching Point}

A species cannot be placed under unrelated genera without explanation.

\[
\boxed{
\text{Genuine classification requires ordered genera.}
}
\]

\section*{14. Higher Genera Must Also Be Predicated in Essence}

If a proposed genus is truly the genus of a species, then its higher genera must
also be predicated of the species in the category of essence.

\[
S \subset G \subset H \subset I
\]

\[
G,H,I \text{ must all be predicated essentially of } S
\]

\subsection*{Example}

If animal is the genus of man, then living thing and substance must also be
predicated essentially of man.

\[
\text{man} \subset \text{animal} \subset \text{living thing} \subset \text{substance}
\]

\subsection*{Teaching Point}

A true genus brings its higher genera with it.

\[
\boxed{
\text{Look upward through the hierarchy to test the proposed genus.}
}
\]

\section*{15. If a Higher Genus Fails, the Lower Genus Fails}

If any higher genus of the proposed genus fails to be predicated essentially of
the species, the proposed genus is false.

\subsection*{Pattern}

\[
S \subset G \subset H
\]

If:

\[
H \text{ is not essentially predicated of } S
\]

then:

\[
G \text{ is not the genus of } S
\]

\subsection*{Teaching Point}

The chain of essence cannot break above the proposed genus.

\[
\boxed{
\text{A false higher genus exposes a false lower genus.}
}
\]

\section*{16. Test the Things of Which the Species Is Predicated}

If the given species is itself predicated of other things, then the proposed
genus must also be predicated essentially of those things.

\subsection*{Pattern}

If:

\[
A \subset S
\]

and:

\[
S \subset G
\]

then:

\[
A \subset G
\]

\subsection*{Teaching Point}

Whatever falls under the species should also fall under the species' genus.

\[
\boxed{
\text{The genus follows the species down to its instances.}
}
\]

\section*{17. Definitions of Genera Must Apply to the Species}

The definition of the genus must apply to the species and to the things that
partake of the species.

\[
\text{definition of } G \text{ applies to } S
\]

\[
\text{definition of } G \text{ applies to instances of } S
\]

\subsection*{Teaching Point}

It is not enough for the name of the genus to apply; its definition must also
apply.

\[
\boxed{
\text{A true genus is predicated synonymously, with its definition.}
}
\]

\section*{18. Do Not Render Differentia as Genus}

A differentia should not be treated as a genus.

\subsection*{Example}

Do not call immortal the genus of God.

\[
\text{God} = \text{living being} + \text{immortal}
\]

Here:

\[
\text{genus} = \text{living being}
\]

\[
\text{differentia} = \text{immortal}
\]

Immortal tells what sort of living being, not the broad kind itself.

\subsection*{Teaching Point}

The genus tells the broad kind; the differentia tells what sort within that
kind.

\[
\boxed{
\text{Differentia qualifies the genus; it is not itself the genus.}
}
\]

\section*{19. Do Not Place the Differentia Inside the Genus as Species}

A differentia is not a species of the genus it divides.

\subsection*{Example}

Odd is not a species of number.

\[
\text{odd} \neq \text{a number}
\]

Rather:

\[
\text{odd} = \text{differentia of number}
\]

Odd divides numbers into a kind.

\subsection*{Teaching Point}

A differentia divides a genus; it is not itself a species under that genus.

\[
\boxed{
\text{Differentiae are not species or individuals.}
}
\]

\section*{20. Do Not Place Genus Inside Species}

Sometimes a person places the genus inside the species by mistake.

\subsection*{Examples}

\[
\text{contact} \neq \text{juncture}
\]

Rather:

\[
\text{juncture} \Rightarrow \text{contact}
\]

but not:

\[
\text{contact} \Rightarrow \text{juncture}
\]

Likewise:

\[
\text{mixture} \neq \text{fusion}
\]

and:

\[
\text{locomotion} \neq \text{carriage}
\]

\subsection*{Teaching Point}

If the supposed species is wider than the supposed genus, the classification is
reversed.

\[
\boxed{
\text{Misordered extension exposes false genus claims.}
}
\]

\section*{21. Do Not Place Species Inside Differentia}

The species should not be treated as if it were inside the differentia.

\subsection*{Example}

Do not say:

\[
\text{immortal} = \text{a god}
\]

Immortal has wider denotation than God, because other living beings might be
called immortal in some context.

\subsection*{Teaching Point}

The differentia is at least as wide as the species it marks.

\[
\boxed{
\text{Species is not the genus of its own differentia.}
}
\]

\section*{22. Do Not Place Genus Inside Differentia}

The genus should not be placed inside the differentia.

\subsection*{Examples}

Do not define colour as what pierces the vision if piercing is a differentia.

Do not define number as what is odd.

\[
\text{number} \neq \text{odd}
\]

Odd is a way of dividing number, not the genus of number.

\subsection*{Teaching Point}

The genus has wider denotation than its differentia and does not partake of its
differentia.

\[
\boxed{
\text{The genus is not contained inside one of its divisions.}
}
\]

\section*{23. Genus and Differentia Must Accompany the Species}

If a proposed genus or differentia can be absent from the alleged species while
the species remains, it is not really the genus or differentia.

\[
S \text{ exists} \Rightarrow G \text{ present}
\]

\[
S \text{ exists} \Rightarrow D \text{ present}
\]

\subsection*{Example}

If movement can be absent from the soul, then movement is not the genus of soul.

If truth and falsehood can be absent from opinion, then they are not its
differentia.

\subsection*{Teaching Point}

The genus and differentia do not come and go while the species remains.

\[
\boxed{
\text{Essential constituents accompany the species as long as it exists.}
}
\]

\section*{24. The Species Must Not Partake of a Contrary Genus}

If the species can partake of a contrary of the proposed genus, the proposed
genus fails.

\subsection*{Pattern}

\[
S \subset G
\]

\[
S \subset \text{contrary of } G
\]

\[
\therefore G \text{ is not the genus of } S
\]

\subsection*{Teaching Point}

A thing cannot essentially belong under contrary genera at the same time.

\[
\boxed{
\text{Contrary genera expose false classification.}
}
\]

\section*{25. The Species Must Not Have What the Genus Cannot Have}

If the species has a feature impossible for any member of the proposed genus,
then the proposed genus is false.

\subsection*{Example}

If the soul shares in life, while no number can live, then soul cannot be a
species of number.

\[
\text{soul has life}
\]

\[
\text{number cannot have life}
\]

\[
\therefore \text{soul is not a number}
\]

\subsection*{Teaching Point}

A proposed genus fails if the species has features impossible within that genus.

\[
\boxed{
\text{The species must be possible within the nature of the genus.}
}
\]

\section*{26. Genus and Species Must Be Synonymous, Not Homonymous}

A true genus is predicated synonymously of its species.

That means the name and the definition apply in the same sense.

\[
\text{true genus} = \text{synonymous predication}
\]

Mere shared naming is homonymy.

\[
\text{mere shared word} = \text{homonymy}
\]

\subsection*{Teaching Point}

A genus must be predicated in the same sense, not merely with the same word.

\[
\boxed{
\text{Homonymy is not classification.}
}
\]

\section*{27. A Genus Must Have More Than One Species}

Every genus has more than one species.

\[
\text{genus} \Rightarrow \text{multiple species}
\]

If it is impossible for there to be another species besides the one given, then
the proposed term is not really a genus.

\subsection*{Teaching Point}

A genus is a common kind, not a private label for one species.

\[
\boxed{
\text{Where there is no possible division into species, there is no genus.}
}
\]

\section*{28. Do Not Use Metaphor as Genus}

A metaphorical expression should not be used as a genus.

\subsection*{Example}

Do not define temperance as a harmony.

\[
\text{temperance is a harmony}
\]

This may be metaphorically suggestive, but it is not literal classification.

\subsection*{Teaching Point}

A genus must classify literally, not poetically.

\[
\boxed{
\text{Metaphor is not genus.}
}
\]

\section*{29. Use Contraries to Test Genus}

Contraries provide tests for proposed genera.

If the genus has no contrary, then contrary species should be in the same genus.

\[
S \text{ and contrary } S' \subset G
\]

If the genus has a contrary, then the contrary species should be found in the
contrary genus.

\[
S \subset G
\]

\[
S' \subset G'
\]

\subsection*{Example}

If virtue and vice are contrary genera, then justice and injustice should fall
under them respectively.

\[
\text{justice} \subset \text{virtue}
\]

\[
\text{injustice} \subset \text{vice}
\]

\subsection*{Teaching Point}

Contrary species should be located according to the contrary structure of their
genera.

\[
\boxed{
\text{Contrary structure should be preserved across genus and species.}
}
\]

\section*{30. Intermediates Must Fit the Same Structure}

If contraries have intermediates, their genera should have intermediates in a
corresponding way.

Likewise, if the genera have intermediates, the species should have
intermediates in a corresponding way.

\subsection*{Example}

Virtue and vice have intermediates.

Justice and injustice have intermediates.

The structures correspond.

\subsection*{Teaching Point}

The structure of extremes and intermediates should be preserved across genus
and species.

\[
\boxed{
\text{A proposed genus should preserve the opposition-structure of the species.}
}
\]

\section*{31. Co-Ordinates and Inflected Forms}

Aristotle recommends looking at co-ordinates and inflected forms.

\subsection*{Example}

If justice is a kind of knowledge, then:

\[
\text{justly} = \text{knowingly}
\]

and:

\[
\text{the just man} = \text{a man of knowledge}
\]

If any of these related expressions fails, the genus claim is weakened.

\subsection*{Teaching Point}

A genus claim should hold across the family of related expressions.

\[
\boxed{
\text{Test the genus through its linguistic relatives.}
}
\]

\section*{32. Likenesses and Analogies}

Aristotle also uses likenesses.

\subsection*{Example}

\[
\text{pleasant}:\text{pleasure}
=
\text{useful}:\text{good}
\]

If pleasure is a kind of good, then the pleasant may be treated as a kind of
useful.

\subsection*{Teaching Point}

Analogous relations can transfer pressure onto a proposed genus.

\[
\boxed{
\text{Analogy gives a route for testing genus, but not an automatic proof.}
}
\]

\section*{33. Generation, Destruction, Capacities, and Uses}

The same analogical testing applies to:

\[
\begin{array}{ll}
1. & \text{processes of generation} \\
2. & \text{processes of destruction} \\
3. & \text{things that generate} \\
4. & \text{things that destroy} \\
5. & \text{capacities} \\
6. & \text{uses}
\end{array}
\]

\subsection*{Examples}

If to build is to be active, then to have built is to have been active.

If learning is recollection, then having learned is having recollected.

If a capacity is a disposition, then to be capable is to be disposed.

\subsection*{Teaching Point}

A genus claim should remain stable across related processes, products,
capacities, and uses.

\[
\boxed{
\text{The same classification should survive across the term's natural family.}
}
\]

\section*{34. Privation and Possession}

If the opposite of the species is a privation, one can test whether that
privation falls under the opposite genus.

\subsection*{Example}

If blindness is a form of insensibility, then sight is a form of sensation.

\[
\text{blindness} \subset \text{insensibility}
\]

\[
\therefore \text{sight} \subset \text{sensation}
\]

\subsection*{Teaching Point}

Privation and possession can test whether a genus relation is properly paired.

\[
\boxed{
\text{The positive and privative forms should belong under corresponding genera.}
}
\]

\section*{35. Use Negation by Converting the Order}

If the species falls under the genus, then whatever is not the genus is not the
species.

\[
S \subset G
\Rightarrow
\neg G \Rightarrow \neg S
\]

\subsection*{Example}

If the pleasant is a kind of good, then what is not good is not pleasant.

\[
\text{pleasant} \subset \text{good}
\]

\[
\text{not good} \Rightarrow \text{not pleasant}
\]

\subsection*{Teaching Point}

If the genus is denied, the species must be denied.

\[
\boxed{
\text{A species cannot appear where its genus is absent.}
}
\]

\section*{36. Relative Terms and Their Genera}

If the species is a relative term, the genus is usually expected to be relative
as well.

\subsection*{Example}

\[
\text{double} \subset \text{multiple}
\]

Both are relative terms.

\[
\text{double of something}
\]

\[
\text{multiple of something}
\]

\subsection*{Teaching Point}

Relative species should fit the relational grammar of their genera.

\[
\boxed{
\text{Relatives bring their correlates with them.}
}
\]

\section*{37. Same Relational Use}

If the species is used in relation to something, the genus should be used in a
corresponding relation.

\subsection*{Example}

If double means double of a half, then multiple should also be used in relation
to that same correlate.

\[
\text{double of } X
\Rightarrow
\text{multiple of } X
\]

\subsection*{Teaching Point}

A proposed genus should preserve the relational pattern of its species.

\[
\boxed{
\text{Genus and species should have compatible relational grammar.}
}
\]

\section*{38. Case-Relationships and Constructions}

Aristotle also examines grammatical constructions.

If genus and species are truly related, their constructions should often match.

\subsection*{Examples}

We say:

\[
\text{double of something}
\]

and:

\[
\text{multiple of something}
\]

We say:

\[
\text{knowledge of something}
\]

and:

\[
\text{object of knowledge}
\]

\subsection*{Teaching Point}

The grammar of a term can reveal whether it belongs under the proposed genus.

\[
\boxed{
\text{Logical relations often show themselves in linguistic constructions.}
}
\]

\section*{39. Different Kinds of Relatives}

Some relative terms must be found in the things to which they are related.

Others may be found there but need not be.

Others cannot be found there at all.

\subsection*{Examples}

A state or disposition is found in the thing whose state or disposition it is.

An object of knowledge may be known by a soul, but need not itself contain the
knowledge.

A contrary is not found in its contrary.

\subsection*{Teaching Point}

Do not place a relative term under a genus that has the wrong relational mode.

\[
\boxed{
\text{Relatives must be classified according to how they relate to their correlates.}
}
\]

\section*{40. State and Activity}

Aristotle warns against placing a state in the genus of activity, or an activity
inside the genus of state.

\[
\text{state} \neq \text{activity}
\]

\subsection*{Example}

If sensation is a state, it should not be defined as movement through the body.

\[
\text{sensation} \neq \text{movement through the body}
\]

\subsection*{Teaching Point}

A proposed genus fails if it places the thing in the wrong mode of being.

\[
\boxed{
\text{State and activity must not be confused.}
}
\]

\section*{41. Capacity and State}

A state should not be ranked within the capacity that attends it.

\subsection*{Examples}

Do not define good temper as control of anger.

Do not define courage as control of fears.

Do not define justice as control of gains.

\subsection*{Teaching Point}

Virtue is not mere control-capacity.

\[
\boxed{
\text{A stable excellence is not the same as the ability to resist passion.}
}
\]

\section*{42. Virtue and Self-Control}

The courageous or good-tempered person is not merely the person who feels
passion and controls it.

Rather, the virtuous person is properly ordered with respect to those passions.

The self-controlled person still has contrary passions, but resists them.

\[
\text{self-control} \neq \text{virtue}
\]

\[
\text{control of passion} \neq \text{proper ordering of passion}
\]

\subsection*{Teaching Point}

Do not mistake an accompanying capacity for the genus of a virtue.

\[
\boxed{
\text{Virtue is deeper than resistance.}
}
\]

\section*{43. Do Not Mistake Attendant Features for Genera}

An attendant feature is not necessarily a genus.

\subsection*{Example: Anger and Pain}

Pain may accompany anger.

But anger is not a species of pain.

\[
\text{anger has pain} \neq \text{anger is pain}
\]

\subsection*{Example: Conviction and Conception}

Conception may accompany conviction.

But conviction is not a species of conception.

One may have the same conception without being convinced of it.

\subsection*{Teaching Point}

What accompanies a thing need not state what the thing is.

\[
\boxed{
\text{Accompaniment is not essence.}
}
\]

\section*{44. Genus and Species Must Occur in the Same Subject}

If a genus truly contains a species, then what contains the species should also
contain the genus.

\[
\text{contains species} \Rightarrow \text{contains genus}
\]

\subsection*{Example}

If shame is fear, then shame and fear should occur in the same faculty.

But shame is found in the reasoning faculty, while fear is found in the spirited
faculty.

Therefore, shame is not fear.

\subsection*{Example}

If anger is pain, then anger and pain should occur in the same faculty.

But pain is found in the faculty of desire, while anger is found in the spirited
faculty.

Therefore, anger is not pain.

\subsection*{Teaching Point}

A genus and its species must naturally belong in the same subject.

\[
\boxed{
\text{Different natural locations expose false genus claims.}
}
\]

\section*{45. Genus Is Not Imparted Only in a Respect}

A genus is not said of its species merely in a particular respect.

\[
\text{genus belongs simply}
\]

not:

\[
\text{genus belongs only in a respect}
\]

\subsection*{Example}

A man is not an animal merely in respect of his body.

He is animal simply.

\subsection*{False Genus Example}

Animal should not be defined as object of sight.

An animal is visible only in respect of its body, not in respect of its soul.

\subsection*{Teaching Point}

A genus belongs to its species without qualification.

\[
\boxed{
\text{Qualified predication is not generic predication.}
}
\]

\section*{46. Do Not Place Whole Inside Part}

Sometimes people place the whole inside the part.

\subsection*{Example}

Animal is sometimes defined as:

\[
\text{animate body}
\]

But body is a part of animal.

A part is not predicated of the whole in the way a genus is.

\[
\text{body} = \text{part of animal}
\]

\[
\text{part} \neq \text{genus of whole}
\]

\subsection*{Teaching Point}

A part is not the genus of the whole.

\[
\boxed{
\text{Material component and genus are not the same relation.}
}
\]

\section*{47. Capacity Is Not the Genus of Blameworthy Things}

Do not put blameworthy things into the genus of capacity.

\subsection*{Example}

A thief is not simply:

\[
\text{one capable of stealing}
\]

Even a good man may be capable of doing evil.

Bad people are named bad because of choice, not mere capacity.

\[
\text{bad character} \rightarrow \text{choice}
\]

\[
\text{not mere capacity}
\]

\subsection*{Teaching Point}

Blame attaches to choice, not bare capacity.

\[
\boxed{
\text{Capacity can be neutral or desirable; vice involves choice.}
}
\]

\section*{48. Intrinsic Goods Are Not Mere Capacities}

Things precious or desirable for their own sake should not be placed in the
genus of capacity, capability, or productivity.

Capacity and productivity are desirable for the sake of something else.

\[
\text{capacity} = \text{for the sake of another}
\]

\[
\text{precious for itself} \neq \text{mere capacity}
\]

\subsection*{Teaching Point}

Intrinsic goods should not be reduced to instrumental powers.

\[
\boxed{
\text{What is desirable for itself is not merely productive of something else.}
}
\]

\section*{49. Some Things Require More Than One Genus}

Some things cannot be placed in a single genus.

\subsection*{Example}

A cheat or slanderer requires both will and capacity.

\[
\text{cheat} = \text{will} + \text{capacity}
\]

A person with will but no capacity is not yet a cheat.

A person with capacity but no will is not yet a cheat.

\subsection*{Teaching Point}

Some classifications require combined genera rather than a single genus.

\[
\boxed{
\text{Complex characters may require more than one explanatory source.}
}
\]

\section*{50. Do Not Reverse Genus and Differentia}

Sometimes people render genus and differentia in converse order.

\subsection*{Examples}

\[
\text{astonishment} = \text{excess of wonderment}
\]

\[
\text{conviction} = \text{vehemence of conception}
\]

Here, wonderment and conception are the genera.

Excess and vehemence are differentiae.

\[
\text{genus} = \text{wonderment}
\]

\[
\text{differentia} = \text{excessive}
\]

\subsection*{Teaching Point}

Do not turn the modifier into the genus.

\[
\boxed{
\text{The base kind is genus; the modifying intensity is differentia.}
}
\]

\section*{51. Absurd Results of Reversal}

If excess were the genus of astonishment, then astonishment would be found in
excess.

This would lead to strange expressions such as:

\[
\text{excessive excess}
\]

If vehemence were the genus of conviction, then conviction would be found in
vehemence, leading to:

\[
\text{vehement vehemence}
\]

\subsection*{Teaching Point}

A false genus often produces absurd predications.

\[
\boxed{
\text{Test a genus by the consequences of its wording.}
}
\]

\section*{52. Affection Is Not the Affected Object}

Do not put an affection into the affected object as its genus.

\subsection*{Example}

Immortality should not be defined simply as everlasting life if immortality is
an affection or feature of life.

\[
\text{immortality} \neq \text{life}
\]

Rather:

\[
\text{immortality} = \text{a feature of life}
\]

\subsection*{Teaching Point}

An affection of a thing is not the genus of the thing affected.

\[
\boxed{
\text{Do not confuse a feature with the subject that bears it.}
}
\]

\section*{53. Movement Is Not the Moved Subject}

Do not define an affection by the object of which it is an affection.

\subsection*{Example}

Wind should not be defined simply as:

\[
\text{air in motion}
\]

Rather, wind is better understood as:

\[
\text{movement of air}
\]

The same air persists whether moving or still.

So wind is not simply air.

\subsection*{Other Examples}

Snow is not simply water.

Mud is not simply earth.

Wine is not simply fermented water.

\subsection*{Teaching Point}

Do not confuse movement or alteration with the subject moved or altered.

\[
\boxed{
\text{The affected subject is not automatically the genus of the affection.}
}
\]

\section*{54. A Term That Is Genus of Nothing Cannot Be This Genus}

If a proposed term is not the genus of anything at all, it cannot be the genus
of the species in question.

\subsection*{Example}

White is not the genus of anything if white objects do not differ specifically
from one another as species of white.

\subsection*{Teaching Point}

A genus must be capable of containing species that differ in kind.

\[
\boxed{
\text{If there are no species under it, it is not a genus.}
}
\]

\section*{55. Beware Universal Predicates as Genera}

Do not name as genus or differentia some feature that follows everything.

\subsection*{Examples}

\[
\text{Being}
\]

\[
\text{Unity}
\]

These follow everything.

If Being were a genus, then everything would be a species of Being.

This collapses the genus-species relation.

\subsection*{Teaching Point}

A predicate too universal to divide into species is not a usable genus.

\[
\boxed{
\text{What follows everything cannot function as an ordinary genus.}
}
\]

\section*{56. Genus Is True Of, Not Merely Inherent In}

The genus is said to be true of the species.

An accident is said to be inherent in the subject.

\subsection*{Example}

White is inherent in snow.

But white is not true of snow as genus.

\[
\text{white inherent in snow} \neq \text{white genus of snow}
\]

\subsection*{Teaching Point}

The grammar of genus differs from the grammar of accident.

\[
\boxed{
\text{Genus is predicated of what the species is, not merely found in it.}
}
\]

\section*{57. Better Species Should Not Be Placed in Worse Genera}

If species and genus both have contraries, beware placing the better species
inside the worse genus.

\subsection*{Pattern}

\[
\text{better species} \rightarrow \text{better genus}
\]

\[
\text{worse species} \rightarrow \text{worse genus}
\]

\subsection*{Example}

If soul is defined as a form of motion, this may be suspect because soul is also
a principle of rest.

If rest is better than motion, then soul should not be placed under the worse
side.

\subsection*{Teaching Point}

The rank of contraries should be preserved across genus and species.

\[
\boxed{
\text{Better things should not be classified under worse genera.}
}
\]

\section*{58. Greater and Less}

Use greater and less degrees to test a proposed genus.

If the proposed genus admits more and less, the species and terms derived from
it should also admit more and less.

\[
G \text{ admits more/less}
\Rightarrow
S \text{ should admit more/less}
\]

\subsection*{Example}

If virtue admits greater and lesser degrees, then justice and the just man
should also admit greater and lesser degrees.

\[
\text{more virtuous}
\]

\[
\text{more just}
\]

\subsection*{Teaching Point}

Degree behavior should be compatible between genus and species.

\[
\boxed{
\text{If the genus varies by degree, its species should fit that pattern.}
}
\]

\section*{59. More Likely Genera}

When several predicates seem to be predicated in the category of essence, and it
is unclear which is the genus, compare the more likely candidates.

If the more likely candidate is not the genus, then the less likely one is not
the genus either.

\[
\text{more likely fails} \Rightarrow \text{less likely fails}
\]

If the less likely candidate is a genus, then the more likely one is also a
genus.

\[
\text{less likely succeeds} \Rightarrow \text{more likely succeeds}
\]

\subsection*{Teaching Point}

Comparative plausibility can test uncertain genus claims.

\[
\boxed{
\text{More-and-less reasoning applies even to classification.}
}
\]

\section*{60. More Likely Species}

The same method applies to species.

If two things have a like claim to be species of a genus, then if one is a
species, so is the other.

If the less likely one is a species, the more likely one is also a species.

\subsection*{Teaching Point}

Comparable cases should be treated consistently.

\[
\boxed{
\text{If the weaker case qualifies, the stronger case qualifies as well.}
}
\]

\section*{61. Establishing a Genus by Multiple Species}

To establish a genus, see whether the genus is predicated in the category of
essence of several different species.

If it is, then it is clearly functioning as a genus.

\subsection*{Pattern}

\[
S_1 \subset G
\]

\[
S_2 \subset G
\]

\[
S_3 \subset G
\]

\[
\therefore G \text{ is a genus}
\]

\subsection*{Teaching Point}

A genus shows itself by essential predication across multiple species.

\[
\boxed{
\text{The genus is common essence across different species.}
}
\]

\section*{62. Distinguishing Genus from Differentia}

Aristotle gives direct criteria for distinguishing genus and differentia.

\[
\begin{array}{ll}
1. & \text{The genus has wider denotation than the differentia.} \\
2. & \text{In stating essence, genus is more fitting than differentia.} \\
3. & \text{Differentia signifies a quality of the genus.}
\end{array}
\]

\subsection*{Example}

Calling man an animal better states what man is than calling him walking.

\[
\text{animal} = \text{genus}
\]

\[
\text{walking} = \text{quality or differentia-like feature}
\]

\subsection*{Teaching Point}

Genus names the kind; differentia qualifies the kind.

\[
\boxed{
\text{Genus tells what it is; differentia tells what sort of that kind it is.}
}
\]

\section*{63. Establishing Genus Through Activity}

If a term is a kind of activity or state, examine whether the subject, in doing
the relevant action, is doing what belongs to the proposed genus.

\subsection*{Example}

To show that knowledge is a form of conviction, ask whether the knower, in
knowing, is convinced.

\[
\text{knowing} \Rightarrow \text{being convinced}
\]

If so, knowledge may be treated as a kind of conviction.

\subsection*{Teaching Point}

Sometimes genus is tested by what the subject is doing in the act named.

\[
\boxed{
\text{Examine the activity to test the classification.}
}
\]

\section*{64. What Always Follows Is Not Necessarily Genus}

Something may always follow a thing without being its genus.

\[
G \text{ genus of } S
\Rightarrow
G \text{ always follows } S
\]

But:

\[
G \text{ always follows } S
\not\Rightarrow
G \text{ genus of } S
\]

\subsection*{Examples}

Rest always follows calm.

Divisibility follows number.

But not every constant follower is a genus.

\subsection*{Teaching Point}

Constant accompaniment is not the same as genus.

\[
\boxed{
\text{What always follows may still fail to say what the thing is.}
}
\]

\section*{65. The Case of Coming-to-Be and Not-Being}

Aristotle gives a sharp example.

Not-being always follows what is coming to be, because what is coming to be is
not yet.

But not-being is not the genus of coming-to-be.

Why?

Because not-being has no species.

\[
\text{not-being follows coming-to-be}
\]

but:

\[
\text{not-being has no species}
\]

\[
\therefore \text{not-being is not a genus}
\]

\subsection*{Teaching Point}

A constant follower must still satisfy the requirements of genus.

\[
\boxed{
\text{Accompaniment alone cannot overcome failure to divide into species.}
}
\]

\section*{66. The Architecture of Book IV}

Book IV moves through the following arc:

\[
\text{what genus is}
\rightarrow
\text{genus vs. accident}
\rightarrow
\text{same category}
\rightarrow
\text{wider denotation}
\]

\[
\rightarrow
\text{species must fit divisions}
\rightarrow
\text{higher genera}
\rightarrow
\text{genus vs. differentia}
\rightarrow
\text{contraries and intermediates}
\]

\[
\rightarrow
\text{co-ordinates and likeness}
\rightarrow
\text{relative terms}
\rightarrow
\text{state, activity, capacity}
\rightarrow
\text{attendant features}
\]

\[
\rightarrow
\text{whole and part}
\rightarrow
\text{universal predicates}
\rightarrow
\text{greater and less}
\rightarrow
\text{constant accompaniment vs. genus}
\]

\subsection*{Teaching Point}

Book IV is Aristotle's guide to avoiding false classification.

\[
\boxed{
\text{The central question is whether the proposed genus really names what the thing is broadly.}
}
\]

\section*{67. Summary of Main Ideas}

\begin{enumerate}[itemsep=0.5em]
    \item Book IV gives commonplace rules for testing proposed genera.

    \item A genus tells what something is broadly.

    \item A true genus is predicated in the category of essence.

    \item A proposed genus must apply to every member of the species.

    \item A genus must not be confused with an accident.

    \item Genus and species must belong to the same category.

    \item The species partakes of the genus, but the genus does not partake of
    the species.

    \item The genus must be wider than the species.

    \item If the genus applies, some species of that genus must apply.

    \item Things not specifically different should have the same genus.

    \item If a species falls under multiple genera, those genera must be
    ordered.

    \item A true genus brings its higher genera with it.

    \item A differentia must not be treated as a genus.

    \item A differentia divides a genus but is not itself a species.

    \item The genus must not be placed inside the species or inside the
    differentia.

    \item The genus and differentia must accompany the species as long as it
    exists.

    \item A species cannot belong under contrary genera.

    \item A species cannot have features impossible for the proposed genus.

    \item Genus and species must be synonymous, not homonymous.

    \item A genus must have more than one species.

    \item Metaphor is not genus.

    \item Contraries and intermediates help test whether a proposed genus is
    correct.

    \item Co-ordinates, inflected forms, likenesses, and analogies can test genus
    claims.

    \item Relative species should fit the relational grammar of their genera.

    \item State, activity, capacity, and disposition must not be confused.

    \item Virtue should not be reduced to mere control-capacity.

    \item Attendant features are not genera.

    \item Genus and species must occur in the same natural subject.

    \item A genus belongs simply, not merely in a certain respect.

    \item A part is not the genus of the whole.

    \item Blameworthy things are not defined by bare capacity.

    \item Intrinsic goods should not be reduced to instrumental powers.

    \item Some complex characters require more than one genus.

    \item Universal predicates such as Being and Unity cannot function as
    ordinary genera.

    \item Greater-and-less reasoning can test genus claims.

    \item Constant accompaniment is not the same as genus.
\end{enumerate}

\section*{Final Framing}

\[
\boxed{
\textit{Topics}, \text{ Book IV, teaches how to test whether a proposed genus really names what a thing is broadly.}
}
\]

Book IV is Aristotle's guide to avoiding false classification. A genus must be
wider than the species, predicated in essence, in the same category, literally
and synonymously said, accompanied by its higher genera, and not confused with
accident, differentia, capacity, activity, metaphor, part, or mere accompaniment.

\[
\boxed{
\text{A genus is not what merely accompanies a thing, but what it is broadly.}
}
\]

\end{document}